\documentclass[11pt,a4paper]{article}

\usepackage[margin=2.5cm,headheight=14pt]{geometry}
\usepackage{hyperref}
\usepackage{titlesec}
\usepackage{enumitem}
\usepackage{booktabs}
\usepackage{array}
\usepackage[table]{xcolor}        % [table] required for \rowcolor
\usepackage{fancyhdr}
\usepackage{mdframed}
\usepackage{listings}
\usepackage[T1]{fontenc}          % T1 encoding required
\usepackage[utf8]{inputenc}       % UTF-8 source encoding
\usepackage{parskip}

\definecolor{codeblue}{RGB}{0,90,160}
\definecolor{sectioncolor}{RGB}{30,60,120}
\definecolor{boxbg}{RGB}{240,245,255}
\definecolor{boxborder}{RGB}{100,140,200}
\definecolor{tablehead}{RGB}{220,230,245}

\hypersetup{
  colorlinks=true,
  linkcolor=codeblue,
  urlcolor=codeblue
}

\titleformat{\section}{\large\bfseries\color{sectioncolor}}{}{0em}{}[\titlerule]
\titleformat{\subsection}{\normalsize\bfseries\color{codeblue}}{}{0em}{}
\titleformat{\subsubsection}{\normalsize\bfseries}{}{0em}{}

\lstset{
  basicstyle=\ttfamily\small,
  backgroundcolor=\color{boxbg},
  frame=single,
  framerule=0.4pt,
  rulecolor=\color{boxborder},
  breaklines=true,
  columns=fullflexible
}

\newmdenv[
  backgroundcolor=boxbg,
  linecolor=boxborder,
  linewidth=1pt,
  innerleftmargin=10pt,
  innerrightmargin=10pt,
  innertopmargin=8pt,
  innerbottommargin=8pt
]{infobox}

\pagestyle{fancy}
\fancyhf{}
\rhead{\textcolor{sectioncolor}{\small OS Development Guide}}
\lhead{\textcolor{sectioncolor}{\small Repository Structure}}
\cfoot{\thepage}
\renewcommand{\headrulewidth}{0.4pt}

\begin{document}

\begin{titlepage}
  \centering
  \vspace*{3cm}
  {\Huge\bfseries\color{sectioncolor} Building a Personal Operating System\par}
  \vspace{0.5cm}
  {\large\color{gray} A Complete Guide to Repository Structure\par}
  \vspace{2cm}
  \hrule height 1pt
  \vspace{0.5cm}
  {\normalsize\color{gray}
    Covers: Firmware \textbullet{} Bootloader \textbullet{} Kernel \textbullet{}
    BSP \textbullet{} Drivers \textbullet{} libc \textbullet{} Userspace \textbullet{} GUI
  \par}
  \vspace{0.5cm}
  \hrule height 0.4pt
  \vfill
  {\small\color{gray} March 2026}
\end{titlepage}

\tableofcontents
\newpage

%% ─────────────────────────────────────────────
\section{Overview: The Boot Chain}

When a computer powers on, code executes in a strict sequence before the OS is usable. Each stage is independent and lives in its own repository.

\begin{infobox}
\textbf{Boot order:}\\
Power on
$\;\rightarrow\;$ BIOS/UEFI (firmware-repo)
$\;\rightarrow\;$ GRUB/U-Boot (bootloader-repo)
$\;\rightarrow\;$ Kernel (kernel-repo)
$\;\rightarrow\;$ Init/Userspace (userspace-repo)
$\;\rightarrow\;$ GUI/Apps
\end{infobox}

\subsection*{Repository layers}

\begin{center}
\begin{tabular}{>{\bfseries}p{4cm} p{9cm}}
\toprule
\rowcolor{tablehead} Repository & Responsibility \\
\midrule
firmware-repo   & BIOS/UEFI, POST, FOTA, hardware init \\
bootloader-repo & GRUB/U-Boot, loads kernel into RAM \\
bsp-repo        & Board Support Package, chip-specific config \\
kernel-repo     & Memory, scheduler, syscalls, drivers \\
drivers-repo    & GPU, USB, storage, network, audio \\
libc-repo       & C standard library, syscall wrappers \\
filesystem-repo & VFS, ext4, FAT32 \\
userspace-repo  & init, shell, core utilities \\
gui-repo        & Window manager, display server \\
apps-repo       & Browser, settings, bundled tools \\
\bottomrule
\end{tabular}
\end{center}

%% ─────────────────────────────────────────────
\section{Firmware Repository (\texttt{firmware-repo})}

The firmware repo holds every piece of code that executes before the OS. It lives in ROM or flash on the motherboard.

\subsection{Hardware Initialisation}
The very first code that runs when power is applied. Responsibilities include:
\begin{itemize}[noitemsep]
  \item POST (Power-On Self-Test) --- validates CPU, RAM, storage
  \item CPU microcode initialisation
  \item DDR/LPDDR RAM training (configuring memory timings and channels)
  \item PCIe link setup and clock/voltage configuration
\end{itemize}
Written in C and assembly; runs before any OS concept exists.

\subsection{UEFI / BIOS Core}
Modern firmware follows the UEFI specification with three phases:
\begin{itemize}[noitemsep]
  \item \textbf{PEI} (Pre-EFI Initialisation) --- early hardware setup
  \item \textbf{DXE} (Driver Execution Environment) --- loads firmware drivers
  \item \textbf{BDS} (Boot Device Selection) --- finds and launches the bootloader
\end{itemize}
Also handles Secure Boot key enrollment and UEFI runtime services.

\subsection{FOTA Engine (Firmware Over-The-Air)}
Manages wireless or remote firmware updates. Steps:
\begin{enumerate}[noitemsep]
  \item Download signed update package
  \item Cryptographically verify the package
  \item Write to the inactive flash slot (A/B scheme)
  \item On next boot, activate the new slot; rollback on failure
\end{enumerate}
Critical on Android, embedded, and IoT devices.

\subsection{Flash Layout}
Defines how physical flash is partitioned:
\begin{itemize}[noitemsep]
  \item Firmware image location in ROM/flash map
  \item Non-volatile storage (NVRAM) for settings
  \item A/B update slots for safe over-the-air updates
\end{itemize}

\subsection{Firmware-Level Drivers}
Early drivers running inside UEFI before the OS kernel loads:
\begin{itemize}[noitemsep]
  \item ACPI tables --- describe hardware topology to the OS
  \item SMBIOS --- system info (vendor, model, serial number)
  \item Minimal NVMe, USB, and framebuffer drivers for the UEFI setup menu
\end{itemize}

\subsection{Security Module}
\begin{itemize}[noitemsep]
  \item TPM (Trusted Platform Module) integration
  \item Measured boot --- records each boot stage into TPM PCRs
  \item Cryptographic key management and signature verification
\end{itemize}

%% ─────────────────────────────────────────────
\section{Bootloader Repository (\texttt{bootloader-repo})}

The bootloader receives control from the firmware and is responsible for loading the kernel into RAM.

\subsection{Stage 1 --- SPL (Secondary Program Loader)}
A tiny piece of code (4--32\,KB) that fits in the chip's internal SRAM before external RAM is initialised:
\begin{itemize}[noitemsep]
  \item Initialises DDR/LPDDR RAM
  \item Loads Stage 2 from flash into RAM
  \item On x86: handled inside GRUB's \texttt{boot.img}
  \item On ARM/Android: U-Boot SPL
\end{itemize}

\subsection{Stage 2 --- Main Loader (GRUB / U-Boot Core)}
The full bootloader, running with complete RAM access:
\begin{itemize}[noitemsep]
  \item Reads boot configuration file
  \item Locates kernel image and initramfs on disk
  \item Loads them into RAM
  \item Sets up boot parameters (kernel command line)
  \item Jumps to the kernel entry point
\end{itemize}

\subsection{Boot Configuration}
Human-readable files that control boot behaviour:
\begin{itemize}[noitemsep]
  \item \texttt{grub.cfg} / \texttt{extlinux.conf}
  \item Kernel image path and command-line arguments
  \item Multi-boot menu entries and timeout
\end{itemize}

\subsection{Signature Verifier}
Core of the Secure Boot chain of trust:
\begin{itemize}[noitemsep]
  \item Checks cryptographic signature on the kernel image
  \item Refuses to load unsigned or incorrectly signed kernels
  \item Signing keys and verification logic both live here
\end{itemize}

\subsection{Device Tree Blobs (DTB)}
Used on ARM and embedded platforms:
\begin{itemize}[noitemsep]
  \item \texttt{.dts} source files compiled to \texttt{.dtb} binaries
  \item Describes hardware layout to the kernel (peripherals, addresses, IRQs)
  \item Bootloader selects the correct DTB for the board and passes it to the kernel
  \item x86 uses ACPI from firmware instead
\end{itemize}

\subsection{Recovery / Fastboot Mode}
Fallback mode triggered by key combination or failed boot:
\begin{itemize}[noitemsep]
  \item Android: \texttt{fastboot} --- flash images over USB, A/B switch, wipe data
  \item Desktop Linux: GRUB rescue shell
  \item Enables recovery without physical reflashing
\end{itemize}

%% ─────────────────────────────────────────────
\section{Board Support Package (\texttt{bsp-repo})}

The BSP is the glue between generic OS software and a specific piece of hardware. It captures all board-specific knowledge that cannot be expressed in the generic kernel or bootloader.

\begin{infobox}
\textbf{Key distinction:} The kernel repo contains \textit{generic} drivers that work across many boards. The BSP repo contains everything \textit{specific to one board or chip family}.
\end{infobox}

\subsection{SoC / CPU Configuration}
\begin{itemize}[noitemsep]
  \item Clock tree setup (frequencies for each clock domain)
  \item Power domain configuration (which chip regions can power down)
  \item Pin multiplexing (assigning functions to GPIO pins)
  \item Voltage regulator settings
\end{itemize}

\subsection{Memory Map}
\begin{itemize}[noitemsep]
  \item Precise map of the address space: RAM, ROM, MMIO registers
  \item Regions reserved for secure world (TrustZone/TEE)
  \item Flash partition layout (bootloader, kernel, userdata slots)
\end{itemize}

\subsection{Device Tree Source (\texttt{.dts})}
\begin{itemize}[noitemsep]
  \item Human-readable hardware description language
  \item Describes every peripheral, its address, interrupt lines, clocks
  \item Compiled to \texttt{.dtb} and passed to the kernel by the bootloader
\end{itemize}

\subsection{HAL (Hardware Abstraction Layer)}
Primarily an Android concept:
\begin{itemize}[noitemsep]
  \item Vendor-provided shared libraries for camera, audio, sensors, GPU, biometrics
  \item Implements standardised Android HAL interfaces
  \item Lets Android run on thousands of devices without knowing each chip's internals
\end{itemize}

\subsection{Board-Specific Drivers}
Drivers too specific for the generic kernel:
\begin{itemize}[noitemsep]
  \item PMIC (Power Management IC) driver
  \item Battery gauge and charger
  \item Display panel initialisation sequence
  \item Touchscreen firmware, fingerprint sensor, NFC controller
\end{itemize}

\subsection{Bootloader Configuration}
U-Boot \texttt{defconfig} files for this board:
\begin{itemize}[noitemsep]
  \item DDR memory training parameters
  \item SPL memory layout
  \item Boot source selection and GPIO boot mode pin
\end{itemize}

%% ─────────────────────────────────────────────
\section{Kernel Repository (\texttt{kernel-repo})}

The kernel is the largest and most complex repository. Its single contract with userspace is the syscall table --- approximately 400 functions that programs invoke via a CPU instruction.

\subsection{\texttt{arch/} --- Architecture-Specific Code}
\begin{itemize}[noitemsep]
  \item Separate subdirectories: \texttt{arch/x86/}, \texttt{arch/arm64/}, \texttt{arch/riscv/}, \ldots
  \item Kernel entry point (first instruction after bootloader jump)
  \item CPU initialisation, interrupt descriptor tables (IDT/GIC)
  \item Context-switching assembly, page table setup
\end{itemize}

\subsection{\texttt{mm/} --- Memory Management}
\begin{itemize}[noitemsep]
  \item Physical page allocator (buddy allocator)
  \item Slab/slub allocator for kernel objects (\texttt{kmalloc}/\texttt{kfree})
  \item Virtual memory areas (VMAs) and page fault handling
  \item Memory-mapped files (\texttt{mmap}), swap, huge pages
  \item OOM (Out of Memory) killer
\end{itemize}

\subsection{\texttt{kernel/} --- Core Kernel}
\begin{itemize}[noitemsep]
  \item CFS (Completely Fair Scheduler) and real-time schedulers
  \item \texttt{fork}/\texttt{exec}/\texttt{exit} process lifecycle
  \item Signal delivery, timer infrastructure (\texttt{hrtimer}, \texttt{jiffies})
  \item Syscall dispatch table
  \item RCU (Read-Copy-Update) synchronisation
  \item \texttt{kthread} infrastructure for kernel threads
\end{itemize}

\subsection{\texttt{fs/} --- Filesystems}
\begin{itemize}[noitemsep]
  \item VFS (Virtual File System) abstraction layer
  \item Filesystem implementations: ext4, btrfs, tmpfs
  \item Virtual filesystems: procfs (\texttt{/proc}), sysfs (\texttt{/sys}), devtmpfs (\texttt{/dev})
  \item Every \texttt{open()}, \texttt{read()}, \texttt{write()} syscall flows through VFS
\end{itemize}

\subsection{\texttt{drivers/} --- Device Drivers}
The largest directory by line count. Organised by category:
\begin{itemize}[noitemsep]
  \item \texttt{drivers/gpu/} --- DRM/KMS display, 3D rendering
  \item \texttt{drivers/usb/} --- host controllers, HID, mass storage
  \item \texttt{drivers/block/} --- NVMe, SATA, eMMC
  \item \texttt{drivers/net/} --- Ethernet, WiFi
  \item \texttt{drivers/input/} --- keyboard, mouse, touchscreen
\end{itemize}

\subsection{\texttt{net/} --- Networking Stack}
\begin{itemize}[noitemsep]
  \item Full TCP/IP implementation: \texttt{net/ipv4/}, \texttt{net/ipv6/}
  \item Socket layer, UDP, ICMP, DNS resolution hooks
  \item Packet filtering (netfilter / iptables / nftables)
  \item Network namespaces
\end{itemize}

\subsection{\texttt{ipc/} --- Inter-Process Communication}
\begin{itemize}[noitemsep]
  \item Pipes and FIFOs
  \item Unix domain sockets
  \item POSIX message queues, semaphores, shared memory
\end{itemize}

\subsection{\texttt{security/} --- Security Framework}
\begin{itemize}[noitemsep]
  \item Linux Security Module (LSM) framework
  \item SELinux and AppArmor implementations
  \item Capabilities (fine-grained privilege control)
  \item Every security-sensitive kernel operation calls through LSM hooks
\end{itemize}

%% ─────────────────────────────────────────────
\section{Drivers Repository (\texttt{drivers-repo})}

When kept separate from the kernel, the drivers repo holds loadable kernel modules (LKMs) and out-of-tree drivers. Every driver follows the same lifecycle: register via \texttt{probe()} on device detection, clean up via \texttt{remove()}.

\subsection{GPU / Display}
\begin{itemize}[noitemsep]
  \item DRM (Direct Rendering Manager) for GPU command submission
  \item KMS (Kernel Mode Setting) for display pipeline and resolution
  \item Mesa implements OpenGL/Vulkan in userspace via DRM \texttt{ioctl}
  \item Fallback: simple linear framebuffer (just write pixels to a memory address)
\end{itemize}

\subsection{Storage}
\begin{itemize}[noitemsep]
  \item NVMe (PCIe-attached SSDs), SATA/AHCI, eMMC/UFS (mobile flash)
  \item All register with the kernel's block layer
  \item Block layer provides the I/O scheduler (ordering and merging reads/writes)
\end{itemize}

\subsection{Network}
\begin{itemize}[noitemsep]
  \item Each NIC driver registers a \texttt{net\_device} with the networking stack
  \item Ethernet, \texttt{mac80211} + firmware for WiFi, \texttt{hci\_core} for Bluetooth
  \item Driver handles raw packets; TCP/IP is handled above in \texttt{net/}
\end{itemize}

\subsection{USB}
\begin{itemize}[noitemsep]
  \item Host controller driver: XHCI (USB 3.x), EHCI (USB 2.0)
  \item Class drivers: HID (keyboards, mice), Mass Storage, CDC (networking/serial)
  \item Fully hotplug --- drivers bind and unbind at runtime
\end{itemize}

\subsection{Input}
\begin{itemize}[noitemsep]
  \item All devices register with the \texttt{evdev} input framework
  \item Exposed as \texttt{/dev/input/eventN} files to userspace
  \item Keyboard: scancode to key event translation
  \item Touchscreen: multi-touch contact point reporting
\end{itemize}

\subsection{Audio}
\begin{itemize}[noitemsep]
  \item ALSA (Advanced Linux Sound Architecture) framework
  \item ASoC (Audio System on Chip) for embedded/mobile
  \item Three driver components: codec driver, platform driver, machine driver
  \item HDMI audio handled as part of the display driver
\end{itemize}

\subsection{Platform / Power Management}
\begin{itemize}[noitemsep]
  \item ACPI (x86) or device tree (ARM) hardware description
  \item \texttt{cpufreq}: CPU frequency scaling (DVFS)
  \item Thermal framework: monitor temperatures, throttle on overheat
  \item Runtime PM: individual devices enter low-power states when idle
  \item Suspend/resume infrastructure for system sleep and wake
\end{itemize}

\begin{infobox}
\textbf{Key rule:} A driver should never contain policy --- it should only talk to hardware and expose a clean interface upward. All decisions about \textit{how} to use the hardware belong in the layers above.
\end{infobox}

%% ─────────────────────────────────────────────
\section{C Standard Library (\texttt{libc-repo})}

libc is the bridge between every userspace program and the kernel. Almost every program links against it, often transparently.

\begin{infobox}
\texttt{your program} $\;\rightarrow\;$ calls \texttt{printf()}, \texttt{malloc()}, \texttt{open()}\\
\texttt{libc} $\;\rightarrow\;$ executes \texttt{syscall} instruction\\
\texttt{kernel} $\;\rightarrow\;$ talks to hardware via drivers
\end{infobox}

\subsection{Syscall Wrappers}
\begin{itemize}[noitemsep]
  \item Wraps every kernel syscall in a clean C function
  \item \texttt{open()}, \texttt{read()}, \texttt{write()}, \texttt{close()}, \texttt{mmap()}, \ldots
  \item Marshals arguments, executes \texttt{syscall} instruction, checks for errors
  \item Without libc you would write this assembly yourself for every OS call
\end{itemize}

\subsection{Memory Allocator}
\begin{itemize}[noitemsep]
  \item \texttt{malloc()}, \texttt{free()}, \texttt{calloc()}, \texttt{realloc()}
  \item Requests large chunks from kernel via \texttt{brk()} / \texttt{mmap()}
  \item Carves those chunks into small allocations for programs
  \item Tracks freed memory and reuses it
  \item Implementations: ptmalloc (glibc), jemalloc, tcmalloc
\end{itemize}

\subsection{stdio --- Buffered I/O}
\begin{itemize}[noitemsep]
  \item \texttt{printf()}, \texttt{scanf()}, \texttt{fopen()}, \texttt{fread()}, \texttt{fwrite()}
  \item Adds a buffering layer over raw \texttt{read()}/\texttt{write()} syscalls
  \item Buffer flushes to kernel in one syscall when full or on \texttt{fflush()}
  \item Makes I/O dramatically faster by batching many small writes
\end{itemize}

\subsection{String and Math}
\begin{itemize}[noitemsep]
  \item \texttt{strlen}, \texttt{memcpy}, \texttt{memmove}, \texttt{strcpy}, \texttt{strcmp}, \texttt{strcat}
  \item Often hand-optimised in SIMD assembly for maximum speed
  \item \texttt{sin}, \texttt{cos}, \texttt{sqrt}, \texttt{pow}, \texttt{log} --- IEEE 754 floating-point
\end{itemize}

\subsection{Process Control}
\begin{itemize}[noitemsep]
  \item \texttt{fork()} --- creates a copy of the current process
  \item \texttt{execve()} --- replaces current process image with a new program
  \item \texttt{exit()} --- terminates with cleanup (flush stdio, run \texttt{atexit()} handlers)
  \item \texttt{wait()}/\texttt{waitpid()} --- parent collects child exit status
  \item \texttt{getenv()}/\texttt{setenv()} --- environment variable management
\end{itemize}

\subsection{Threading (pthreads)}
\begin{itemize}[noitemsep]
  \item \texttt{pthread\_create()} --- spawns a thread via \texttt{clone()} syscall
  \item \texttt{pthread\_mutex\_lock/unlock} --- mutual exclusion
  \item \texttt{pthread\_cond\_wait/signal} --- condition variables
  \item \texttt{pthread\_join()} --- wait for thread completion
  \item Thread-Local Storage (TLS) --- per-thread \texttt{errno} and variables
\end{itemize}

\subsection{Implementations}

\begin{center}
\begin{tabular}{>{\bfseries}p{2.5cm} p{4cm} p{5cm}}
\toprule
\rowcolor{tablehead} Name & Use case & Notes \\
\midrule
glibc   & Linux desktop / server & Full-featured, large, very mature \\
musl    & Embedded / Alpine Linux & Lightweight, clean, \textasciitilde{}30k lines \\
newlib  & Bare-metal / RTOS       & No OS required \\
relibc  & Redox OS / Rust OS      & Written in Rust, C-ABI compatible \\
\bottomrule
\end{tabular}
\end{center}

%% ─────────────────────────────────────────────
\section{Userspace Repository (\texttt{userspace-repo})}

The userspace repo is where your OS identity lives. The kernel is largely the same across many OSes --- what makes each OS \textit{different} is almost entirely the userspace.

\subsection{Init System (PID 1)}
\begin{itemize}[noitemsep]
  \item First userspace process the kernel starts after boot
  \item Parent of all other processes
  \item Mounts filesystems, starts services, sets up networking
  \item Options: systemd (most modern distros), OpenRC (Gentoo, Alpine), runit (Void Linux), or a custom init
  \item Writing a simple init is achievable: fork and exec other programs from C
\end{itemize}

\subsection{Shell}
\begin{itemize}[noitemsep]
  \item Command-line interface between user and kernel
  \item Parses commands, handles pipes (\texttt{|}), redirection (\texttt{>}), environment variables
  \item Options: bash, sh, zsh, fish, or a custom REPL
  \item Minimal shell project: parse input $\rightarrow$ \texttt{fork()} $\rightarrow$ \texttt{execve()}
\end{itemize}

\subsection{Core Utilities}
\begin{itemize}[noitemsep]
  \item All the small programs users and scripts rely on
  \item \texttt{ls}, \texttt{cp}, \texttt{mv}, \texttt{cat}, \texttt{echo}, \texttt{grep}, \texttt{mount}, \texttt{ps}, \texttt{kill}
  \item BusyBox: all utilities in one binary (ideal for embedded/minimal systems)
  \item GNU coreutils: full individual implementations
\end{itemize}

\subsection{Service Manager}
\begin{itemize}[noitemsep]
  \item Manages long-running background daemons
  \item Network manager, syslog/journald, udev, cron, SSH server
  \item Integrated into systemd on modern systems; separate on lighter systems
\end{itemize}

\subsection{Package / Update System}
\begin{itemize}[noitemsep]
  \item How software is installed, updated, and removed
  \item apt (Debian/Ubuntu), pacman (Arch), dnf (Fedora)
  \item For embedded/mobile: custom OTA update system
\end{itemize}

\subsection{OS Branding Layer}
What makes it \textit{your} OS:
\begin{itemize}[noitemsep]
  \item Name and version string in \texttt{/etc/os-release}
  \item Default configuration files and shell prompt
  \item Default wallpaper, preinstalled apps, update server URLs
\end{itemize}

%% ─────────────────────────────────────────────
\section{Rust as an Alternative to C}

Rust is a fully viable modern language for OS-level code, from the kernel down to libc.

\subsection{Why Rust for OS Development?}
\begin{itemize}[noitemsep]
  \item Memory safety without a garbage collector
  \item No use-after-free, no buffer overflows, no null pointer dereferences at compile time
  \item Safe concurrency --- data races are a compile error
  \item Same or better performance as C
  \item These are exactly the bug classes causing most kernel CVEs
\end{itemize}

\subsection{Rust Options for Each Layer}

\begin{center}
\begin{tabular}{>{\bfseries}p{3cm} p{4cm} p{5cm}}
\toprule
\rowcolor{tablehead} Layer & Rust approach & Notes \\
\midrule
Kernel           & \texttt{no\_std} + \texttt{core} crate & Linux accepts Rust drivers since v6.1 \\
Bare-metal        & \texttt{core} + \texttt{alloc}          & Zero libc, full language features \\
libc replacement  & relibc                         & Full C-ABI compatible, used by Redox \\
Userspace         & Rust \texttt{std}               & Wraps libc transparently \\
\bottomrule
\end{tabular}
\end{center}

\subsection{Real-World Rust OS Projects}
\begin{itemize}[noitemsep]
  \item \textbf{Redox OS} --- full OS in Rust: microkernel + relibc + userspace
  \item \textbf{Linux kernel} --- Rust drivers accepted since kernel 6.1 (2022)
  \item \textbf{Android} --- new low-level components written in Rust since 2021
  \item \textbf{Tock OS} --- embedded Rust OS for microcontrollers
\end{itemize}

%% ─────────────────────────────────────────────
\section{Realistic Development Roadmap}

\begin{center}
\begin{tabular}{>{\bfseries}p{3cm} p{9.5cm}}
\toprule
\rowcolor{tablehead} Milestone & Goal \\
\midrule
Week 1--2   & Bootloader prints ``Hello World'' in QEMU \\
Month 1     & Kernel boots, handles interrupts, basic memory management \\
Month 3     & Processes, scheduler, system calls working \\
Month 6     & Filesystem, shell, basic userspace programs \\
Year 1+     & Networking, GUI, driver ecosystem \\
\bottomrule
\end{tabular}
\end{center}

\subsection{Recommended Tools}
\begin{itemize}[noitemsep]
  \item \textbf{QEMU} / VirtualBox --- test your OS without rebooting real hardware
  \item \textbf{GDB} --- debug your kernel with \texttt{-s -S} QEMU flags
  \item \textbf{NASM} / GAS --- assemblers for boot code
  \item \textbf{GCC cross-compiler} or \textbf{clang} --- target your architecture
  \item \textbf{make} / cmake --- build system
\end{itemize}

\subsection{Learning Resources}
\begin{itemize}[noitemsep]
  \item \href{https://wiki.osdev.org}{OSDev Wiki} --- the definitive reference for OS development
  \item Linux 0.01 --- Linus's first kernel, only \textasciitilde{}10,000 lines
  \item xv6 --- MIT's teaching OS, clean and well-documented
  \item SerenityOS --- modern hobby OS with full GUI
  \item Redox OS --- written entirely in Rust
  \item \textit{Operating Systems: Three Easy Pieces} (free online)
  \item \textit{Modern Operating Systems} --- Tanenbaum
\end{itemize}

\vfill
\hrule
\vspace{0.3cm}
{\small\color{gray} This document covers the full repository structure for building a personal operating system, from firmware to GUI. Each repository has a single clear owner, one release cycle, and one responsibility.}

\end{document}
